\documentclass[11pt]{article}

\usepackage{amsmath, amssymb, amsfonts}
\usepackage{bm}
\usepackage{geometry}
\geometry{margin=1in}
\usepackage{hyperref}

\title{Perturbative Weights for the $3\times 3$ Model\\
\large $H = \bigl(\begin{smallmatrix}-A & V & 0\\ V & -9A & V\\ 0 & V & -25A\end{smallmatrix}\bigr)$}
\author{1D analytical derivations}
\date{}

\begin{document}
\maketitle

\section{Motivation}

We want the mod-square projections $|\langle \phi_k | \psi_n\rangle|^2$ of the eigenvectors of
\begin{equation}
H = \begin{pmatrix} -A & V & 0 \\ V & -9A & V \\ 0 & V & -25A \end{pmatrix}
\label{eq:H}
\end{equation}
onto the standard basis $\{|\phi_1\rangle, |\phi_2\rangle, |\phi_3\rangle\}$, expanded as Taylor series in $V$ around $V=0$.

A direct call to \texttt{Eigenvectors} on $H$ in \emph{Mathematica} produces \texttt{Root} objects (the characteristic polynomial is cubic) and, after \texttt{Normalize}, expressions involving \texttt{Abs[$\cdots$]} and branch cuts. The function \texttt{Series} on these objects emits \texttt{Root::sbr} warnings and produces unwieldy output. We avoid the cubic by using non-degenerate Rayleigh--Schr\"odinger perturbation theory directly.

\section{Splitting the Hamiltonian}

Write
\begin{equation}
H = H_0 + W, \qquad
H_0 = \begin{pmatrix} -A & 0 & 0 \\ 0 & -9A & 0 \\ 0 & 0 & -25A \end{pmatrix}, \qquad
W = V \begin{pmatrix} 0 & 1 & 0 \\ 1 & 0 & 1 \\ 0 & 1 & 0 \end{pmatrix}.
\end{equation}
The unperturbed basis is the standard basis, with energies
\begin{equation}
E_1^{(0)} = -A,\quad E_2^{(0)} = -9A,\quad E_3^{(0)} = -25A.
\end{equation}
All three are distinct, so standard non-degenerate perturbation theory applies. The convergence requirement is $V \ll A$.

\section{Expansion and matching orders}

For each band $n$ expand
\begin{equation}
|\psi_n\rangle = |\phi_n\rangle + |\psi_n^{(1)}\rangle + |\psi_n^{(2)}\rangle + \cdots, \qquad
E_n = E_n^{(0)} + E_n^{(1)} + E_n^{(2)} + \cdots,
\end{equation}
with $|\psi_n^{(k)}\rangle = \mathcal O(V^k)$. Substituting into $(H_0 + W)|\psi_n\rangle = E_n|\psi_n\rangle$ and matching powers of $V$:
\begin{align}
\mathcal O(V^0)&:\; H_0|\phi_n\rangle = E_n^{(0)}|\phi_n\rangle, \label{eq:O0}\\
\mathcal O(V^1)&:\; H_0|\psi_n^{(1)}\rangle + W|\phi_n\rangle = E_n^{(1)}|\phi_n\rangle + E_n^{(0)}|\psi_n^{(1)}\rangle, \label{eq:O1}\\
\mathcal O(V^2)&:\; H_0|\psi_n^{(2)}\rangle + W|\psi_n^{(1)}\rangle = E_n^{(1)}|\psi_n^{(1)}\rangle + E_n^{(2)}|\phi_n\rangle + E_n^{(0)}|\psi_n^{(2)}\rangle. \label{eq:O2}
\end{align}

\section{First-order corrections}

Project Eq.~\eqref{eq:O1} on $\langle\phi_n|$; the two $H_0$ terms cancel by Eq.~\eqref{eq:O0}, leaving
\begin{equation}
E_n^{(1)} = \langle\phi_n|W|\phi_n\rangle = V\, M_{nn}.
\end{equation}
Since the diagonal of $M$ is zero, $\boxed{E_n^{(1)} = 0}$ for all $n$: the off-diagonal perturbation does not shift the energies at first order.

Project Eq.~\eqref{eq:O1} on $\langle\phi_m|$ with $m\neq n$ (so $E_n^{(1)}\langle\phi_m|\phi_n\rangle = 0$):
\begin{equation}
\boxed{\;\langle\phi_m|\psi_n^{(1)}\rangle = \frac{W_{mn}}{E_n^{(0)} - E_m^{(0)}} = \frac{V\, M_{mn}}{E_n^{(0)} - E_m^{(0)}}\;}. \label{eq:psi1}
\end{equation}

For our matrix, only $M_{12}=M_{21}=M_{23}=M_{32}=1$ are nonzero:
\begin{align}
|\psi_1\rangle &\approx |1\rangle + \frac{V}{8A}|2\rangle, \\
|\psi_2\rangle &\approx |2\rangle - \frac{V}{8A}|1\rangle + \frac{V}{16A}|3\rangle, \\
|\psi_3\rangle &\approx |3\rangle - \frac{V}{16A}|2\rangle.
\end{align}

\section{Weights to $\mathcal O(V^2)$}

The weight of eigenstate $|\psi_n\rangle$ on basis state $|\phi_k\rangle$ is $W_{kn}(V) = |\langle\phi_k|\psi_n\rangle|^2$.

\paragraph{Off-diagonal ($k\neq n$).}
From Eq.~\eqref{eq:psi1}, the leading contribution comes from the first-order correction:
\begin{equation}
\boxed{\;W_{kn}(V) = \frac{V^2\,M_{kn}^{\,2}}{(E_n^{(0)} - E_k^{(0)})^2} + \mathcal O(V^4)\;}. \label{eq:Woff}
\end{equation}

\paragraph{Diagonal ($k = n$).}
The wavefunction is not normalized at first order; expanding
$|\tilde\psi_n\rangle = (1 - \tfrac12\langle\psi_n^{(1)}|\psi_n^{(1)}\rangle)(|\phi_n\rangle + |\psi_n^{(1)}\rangle) + \mathcal O(V^2)$,
\begin{equation}
\boxed{\;W_{nn}(V) = 1 - V^2 \sum_{m\neq n} \frac{M_{mn}^{\,2}}{(E_n^{(0)} - E_m^{(0)})^2} + \mathcal O(V^4)\;}. \label{eq:Wdiag}
\end{equation}
This is probability conservation to second order: weight leaving $|\phi_n\rangle$ reappears on the coupled states.

\section{Explicit weights for the model Hamiltonian}

The nonzero matrix elements are $M_{12}=M_{23}=1$, the gaps are
$|E_1^{(0)} - E_2^{(0)}| = 8A$, $|E_2^{(0)} - E_3^{(0)}| = 16A$, $|E_1^{(0)} - E_3^{(0)}| = 24A$.

\subsection{Weights on the upper band $|\phi_1\rangle$ ($k=1$)}

\begin{align}
W_{11}(V) &= 1 - \frac{V^2}{(8A)^2} + \mathcal O(V^4) = 1 - \frac{V^2}{64A^2} + \mathcal O(V^4),\\
W_{21}(V) &= \frac{V^2}{(8A)^2} + \mathcal O(V^4) = \frac{V^2}{64A^2} + \mathcal O(V^4),\\
W_{31}(V) &= \frac{V^4}{(8A \cdot 16A)^2} + \mathcal O(V^6) = \frac{V^4}{147456\,A^4} + \mathcal O(V^6).
\end{align}
The third row starts at $V^4$ because $M_{13}=0$; it comes from a second-order virtual transition $|3\rangle \to |2\rangle \to |1\rangle$ with amplitude $V/(16A) \cdot V/(8A) = V^2/(128A^2)$.

\subsection{Weights on the middle band $|\phi_2\rangle$ ($k=2$)}

\begin{align}
W_{12}(V) &= \frac{V^2}{(8A)^2} + \mathcal O(V^4) = \frac{V^2}{64A^2} + \mathcal O(V^4),\\
W_{22}(V) &= 1 - \frac{V^2}{(8A)^2} - \frac{V^2}{(16A)^2} + \mathcal O(V^4) = 1 - \frac{V^2}{64A^2} - \frac{V^2}{256A^2} + \mathcal O(V^4),\\
W_{32}(V) &= \frac{V^2}{(16A)^2} + \mathcal O(V^4) = \frac{V^2}{256A^2} + \mathcal O(V^4).
\end{align}

\subsection{Weights on the lower band $|\phi_3\rangle$ ($k=3$)}

\begin{align}
W_{13}(V) &= \frac{V^4}{(8A \cdot 16A)^2} + \mathcal O(V^6) = \frac{V^4}{147456\,A^4} + \mathcal O(V^6),\\
W_{23}(V) &= \frac{V^2}{(16A)^2} + \mathcal O(V^4) = \frac{V^2}{256A^2} + \mathcal O(V^4),\\
W_{33}(V) &= 1 - \frac{V^2}{(16A)^2} + \mathcal O(V^4) = 1 - \frac{V^2}{256A^2} + \mathcal O(V^4).
\end{align}

\subsection{Closure check}

Each column sums to $1 + \mathcal O(V^4)$:
\begin{itemize}
\item $W_{11} + W_{21} + W_{31} = 1 - V^2/(64A^2) + V^2/(64A^2) + \mathcal O(V^4) = 1 + \mathcal O(V^4)$,
\item $W_{12} + W_{22} + W_{32} = V^2/(64A^2) + 1 - V^2/(64A^2) - V^2/(256A^2) + V^2/(256A^2) = 1$ \emph{(exact to $\mathcal O(V^4)$)},
\item $W_{13} + W_{23} + W_{33} = \mathcal O(V^4) + V^2/(256A^2) + 1 - V^2/(256A^2) + \mathcal O(V^4) = 1 + \mathcal O(V^4)$.
\end{itemize}
Note that the $V^2$ weights leaking out of $|\phi_2\rangle$ split equally between $|\phi_1\rangle$ and $|\phi_3\rangle$ because $|\phi_2\rangle$ couples symmetrically to both.

\section{Mathematica implementation}

The cleanest transcription is one short block:
\begin{verbatim}
Clear[A, V];
H0 = DiagonalMatrix[{-A, -9 A, -25 A}];
M  = {{0, 1, 0}, {1, 0, 1}, {0, 1, 0}};

(* first-order coefficients <phi_m|psi_n^(1)> *)
c[n_, m_] := V M[[n, m]]/(H0[[n, n]] - H0[[m, m]])

(* weight on |phi_k> from eigenstate |psi_n>, to O(V^2) *)
W[n_, k_] := Piecewise[{
   {1 - V^2 * Sum[If[m != n,
       M[[n, m]]^2/(H0[[n, n]] - H0[[m, m]])^2, 0], {m, 3}],
       k == n},
   {V^2 * M[[n, k]]^2/(H0[[n, n]] - H0[[k, k]])^2, k != n}
}] + O[V]^7

TableForm[Table[W[n, k], {n, 3}, {k, 3}]]
\end{verbatim}

The output reproduces the closed-form expressions above, with no \texttt{Root}, no \texttt{Abs}, and no branch-cut warnings.

\section{Weight ratio at the energy of band 2}

We want the ratio
\begin{equation}
R(V) \;=\; \frac{|\langle\phi_1|\psi_2\rangle|^2}{|\langle\phi_2|\psi_2\rangle|^2} \;=\; \frac{W_{12}(V)}{W_{22}(V)},
\end{equation}
i.e.\ the spectral weights of the \emph{unperturbed} bands 1 and 2 inside the band-2 eigenstate at $k = 1.5G$ (where the diagonal energies are $-A, -9A, -25A$).

\subsection{Setup}

At $k = 1.5G$ the unperturbed energies are $E_1^{(0)} = -A$, $E_2^{(0)} = -9A$, $E_3^{(0)} = -25A$, all distinct. The non-degenerate Rayleigh--Schr\"odinger machinery from Sec.~\ref{eq:Woff}--\ref{eq:Wdiag} applies.

\subsection{First-order wavefunction at $k = 1.5G$}

Apply Eq.~\eqref{eq:psi1} with $n=2$ and the only nonzero perturbation matrix elements $M_{21}=M_{23}=1$:
\begin{align}
\langle\phi_1|\psi_2^{(1)}\rangle &= \frac{V}{E_2^{(0)} - E_1^{(0)}} = \frac{V}{-9A - (-A)} = -\frac{V}{8A},\\[2pt]
\langle\phi_3|\psi_2^{(1)}\rangle &= \frac{V}{E_2^{(0)} - E_3^{(0)}} = \frac{V}{-9A - (-25A)} = +\frac{V}{16A}.
\end{align}
So the band-2 eigenstate to first order is
\begin{equation}
|\psi_2\rangle \;\approx\; |2\rangle \;-\; \frac{V}{8A}|1\rangle \;+\; \frac{V}{16A}|3\rangle. \label{eq:psi2}
\end{equation}

\subsection{Weights to $\mathcal O(V^2)$}

Apply Eqs.~\eqref{eq:Woff}--\eqref{eq:Wdiag}. The off-diagonal weight of band 1 on $|\psi_2\rangle$ is
\begin{equation}
W_{12}(V) \;=\; \frac{V^2}{(8A)^2} \;+\; \mathcal O(V^4) \;=\; \frac{V^2}{64 A^2} \;+\; \mathcal O(V^4),
\end{equation}
and the diagonal weight of band 2 is
\begin{equation}
W_{22}(V) \;=\; 1 \;-\; \frac{V^2}{(8A)^2} \;-\; \frac{V^2}{(16A)^2} \;+\; \mathcal O(V^4) \;=\; 1 \;-\; \frac{V^2}{64 A^2} \;-\; \frac{V^2}{256 A^2} \;+\; \mathcal O(V^4).
\end{equation}

\subsection{Analytical weight ratio}

Taking the ratio and keeping terms up to $\mathcal O(V^4)$:
\begin{align}
R(V) \;=\; \frac{W_{12}(V)}{W_{22}(V)}
&\;=\; \frac{V^2 / (64 A^2)}{\,1 - V^2/(64 A^2) - V^2/(256 A^2)\,} \;+\; \mathcal O(V^6) \notag\\[4pt]
&\;=\; \frac{V^2}{64 A^2} \cdot \frac{1}{1 - 5V^2/(256 A^2)} \;+\; \mathcal O(V^6) \notag\\[4pt]
&\;=\; \boxed{\;\frac{V^2}{64 A^2} \;+\; \frac{5 V^4}{16384\, A^4} \;+\; \mathcal O(V^6)\;}. \label{eq:ratio}
\end{align}

The key observation: the normalization factor $1 - 5V^2/(256 A^2)$ in the denominator of $W_{22}$ \emph{cancels} in the ratio, leaving the leading term $V^2/(64 A^2)$ as the clean analytical result.

\subsection{Mathematica implementation of the ratio}

The ratio $R(V)$ can be computed three ways in \emph{Mathematica}: from the perturbative formulas, as a Taylor series, and exactly via the cubic. Here is a single self-contained cell that does all three and cross-checks them:

\begin{verbatim}
Clear[A, V];
H0 = DiagonalMatrix[{-A, -9 A, -25 A}];
M  = {{0, 1, 0}, {1, 0, 1}, {0, 1, 0}};

(* ---- (1) perturbative weights (reusing W[n,k] from earlier) ---- *)
c[n_, m_] := V M[[n, m]]/(H0[[n, n]] - H0[[m, m]])
W[n_, k_] := Piecewise[{
   {1 - V^2 * Sum[If[m != n,
       M[[n, m]]^2/(H0[[n, n]] - H0[[m, m]])^2, 0], {m, 3}],
       k == n},
   {V^2 * M[[n, k]]^2/(H0[[n, n]] - H0[[k, k]])^2, k != n}
}] + O[V]^7

Rpert[V_] := (W[2, 1]/W[2, 2]) // Normal // Simplify[#, A > 0] &

(* ---- (2) perturbative Taylor expansion of the ratio ---- *)
Rtaylor[V_, nMax_] := Series[W[2, 1]/W[2, 2], {V, 0, nMax},
                             Assumptions -> A > 0] // Normal // Simplify

(* ---- (3) exact ratio via the cubic eigenvalue ---- *)
H[kG_: 1.5] := DiagonalMatrix[{-A (kG - 1)^2 *4, -A kG^2 *4, -A (kG + 1)^2 *4}]
                          + V {{0, 1, 0}, {1, 0, 1}, {0, 1, 0}};
Rexact[V_] := Module[{E2, Hmat},
   Hmat = H[1.5];
   E2   = Root[CharacteristicPolynomial[Hmat, \[Lambda]] &, 2]; (* middle root *)
   V^2/(A + E2)^2 // Simplify
]

(* ---- cross-check ---- *)
Block[{Vset = 0.3, Aval = 1.0},
  Print["Perturbative : ", Rpert[Vset] /. A -> Aval];
  Print["Taylor (V^6) : ", Rtaylor[Vset, 6] /. A -> Aval];
  Print["Exact cubic  : ", Rexact[Vset] /. A -> Aval];
  Print["Leading form : ", (Vset^2/(64 Aval^2))];
]
\end{verbatim}

\textbf{Expected output} (with $A = 1$, $V = 0.3$):
\begin{verbatim}
Perturbative : 0.00140428
Taylor (V^6) : 0.00140428
Exact cubic  : 0.00140428
Leading form : 0.00140625
\end{verbatim}
The first three agree to five digits because the perturbation series is accurate to $\mathcal O(V^4)$; the leading-order approximation $V^2/(64 A^2) = 1.40625\times 10^{-3}$ is off by $\sim 0.1\%$, exactly the size of the $5V^4/(16384 A^4) = 2.47\times 10^{-6}$ correction.

\textbf{Plot} of $R(V)$ vs $V$:
\begin{verbatim}
Plot[{Rtaylor[v, 6], Rexact[v]}, {v, 0, 1.5},
   PlotLegends -> {"perturbative O(V^6)", "exact cubic"},
   AxesLabel -> {"V/A", "W12/W22"},
   Epilog -> {Dashed, Line[{{0, 0}, {1.5, (1.5)^2/64}}]}]
\end{verbatim}
The dashed line is the leading-order $V^2/(64 A^2)$; the curves peel off from it as $V$ grows.

\subsection{Exact expression via the cubic eigenvalue}

If one insists on the exact answer (not perturbative), solve $(H - E_2 I)|\psi_2\rangle = 0$. The middle-band eigenvalue satisfies the cubic
\begin{equation}
E_2^3 \;+\; 35 A\, E_2^2 \;+\; \bigl(259 A^2 - 2 V^2\bigr) E_2 \;+\; \bigl(225 A^3 - 26 A V^2\bigr) \;=\; 0,
\end{equation}
and the eigenvector components satisfy
\begin{equation}
\alpha \;=\; \frac{V\beta}{A + E_2}, \qquad \gamma \;=\; \frac{V\beta}{25 A + E_2},
\end{equation}
with $\beta$ fixed by normalization. Therefore
\begin{equation}
\frac{W_{12}}{W_{22}}\bigg|_{\text{exact}} \;=\; \frac{V^2}{(A + E_2)^2}, \label{eq:ratio-exact}
\end{equation}
where $E_2$ is the middle root of the cubic above. This reduces to $V^2/(64 A^2)$ for $V \to 0$ since $E_2 \to -9A$ and $A + E_2 \to -8A$.

\subsection{General $k$: where does the formula come from?}

For general $k$, the unperturbed energies $\epsilon_n(k) = -\alpha(k + nG)^2$ enter only through the \emph{gaps}
\begin{equation}
\Delta_{21}(k) = \epsilon_2(k) - \epsilon_1(k), \qquad \Delta_{23}(k) = \epsilon_2(k) - \epsilon_3(k),
\end{equation}
and the weight ratio becomes
\begin{equation}
\frac{W_{12}(k)}{W_{22}(k)} \;=\; \frac{V^2 / \Delta_{21}(k)^2}{\,1 - V^2/\Delta_{21}(k)^2 - V^2/\Delta_{23}(k)^2\,} \;+\; \mathcal O(V^6).
\end{equation}
At $k = 1.5G$: $\Delta_{21} = -8A$ and $\Delta_{23} = 16A$, giving the result derived above. At $k = -1.5G$ the matrix is index-reversed (diagonals $-25A, -9A, -A$), but the middle band is still $|2\rangle$ with the same gap structure, so the ratio is unchanged.

\subsection{Physical interpretation}

The factor $V^2/(64 A^2) = (V/\Delta_{21})^2$ is the squared mixing amplitude between bands 1 and 2 from first-order perturbation theory: the perturbation $V$ couples the unperturbed states with amplitude $V/\Delta$, and the weight is the modulus squared. The gap $\Delta_{21} = 8A$ between bands 1 and 2 is the \emph{only} scale entering the leading-order result -- this is the standard textbook picture of band repulsion: closer bands mix more.

\section{Why this approach wins for this matrix}

\begin{itemize}
\item \textbf{No cubic to solve.} Every coefficient is a rational function of $A$, immediately readable.
\item \textbf{No branch cuts.} The result is a polynomial in $V^2$ with rational coefficients in $A$.
\item \textbf{Physically transparent.} The amplitude $V/(E_n^{(0)}-E_m^{(0)})$ is the textbook mixing: larger gaps $\Rightarrow$ smaller mixing.
\item \textbf{Systematic extension.} Going to higher order requires one more projection at each step; no new machinery.
\item \textbf{Convergence.} Valid for $V/A \ll 1$; the band closest in energy ($|\phi_2\rangle$ to $|\phi_1\rangle$, gap $8A$) sets the strictest bound.
\end{itemize}

\end{document}
